%% kapa: exam=1 round=26 year=2015 period=2 form=A subject=appraisal_law subject_ko=부동산관계법규 seq=20 number=20
%% source: Q-Net 감정평가사 기출문제 (cst00309) artlSeq=5138001 file=extracted/1008/5138001/제26회_감정평가사 1차_기출/제26회_감정평가사 1차_기출/2교시_감평_A형.pdf
%% conversion: kordoc 4.24.51 -> tools/gen1.py (auto); review=proofread-sonnet; flags=none
\begin{question}{20}
부동산가격공시 및 감정평가에 관한 법령상 감정평가법인에 관한 설명으로 옳지 않은 것은?
\begin{choices}
  \item 감정평가법인의 주사무소에 주재하는 최소 감정평가사의 수는 3명이다.
  \item 감정평가법인은 당해 법인의 소속감정평가사 외의 자로 하여금 감정평가업자의 업무를 하게 하여서는 아니된다.
  \item 감정평가법인의 사원 또는 이사는 감정평가사이어야 한다.
  \item 감정평가법인은 사원 전원의 동의 또는 주주총회의 의결이 있는 때에는 국토교통부장관의 인가를 받아 다른 감정평가법인과 합병할 수 있다.
  \item 감정평가법인의 자본금은 2억원 이상이어야 한다.
\end{choices}
\end{question}
\answer{1}
